\scanpage{54}
Alors comme module sur $A_C$, \deleted{$T_p\otimes C$ s'identifie} \inserted{on a un isom} canonique
\[
\unsure{(*)}\qquad\boxed{T_p\otimes_{\mathbb Z_p}C\simeq\prod_\sigma C_\sigma},
\]
où $C_\sigma=C$, comme $C$-vectoriel, l'opération de $A$ dessus se faisant via $\sigma:A\mathrel{\text{---}}C$, où
\[
\sigma\in\operatorname{Hom}_{\mathbb Q_p}(K,C)\xleftarrow{\sim}\operatorname{Aut}_{\mathbb Q_p}(K)\simeq\operatorname{Aut}_{\mathbb Z_p}(A)\simeq\operatorname{Hom}_{\mathbb Z_p}(A,C).
\]
\margin{$*$}
Cette décomposition est compatible avec l'action de $\pi$, de façon précise, $\pi$ opère sur le facteur $C_\sigma$ par $\rho_\sigma(g)$ ($g\in\pi$) donné par
\[
x\in C_\sigma\Longrightarrow\rho_\sigma(g)x=\sigma(\chi_\pi(g))x
\]
Comparons cette décomposition avec celle donnée par le théor. de Tate
\diagram{\(\left\{\begin{gathered}
\begin{tikzcd}
0\arrow[r]&\mathfrak g\otimes C(1)\arrow[r]&T_p\otimes C\arrow[r]\arrow[l,dashed,bend left=35]&\check{\mathfrak g}^{*}\otimes C\arrow[r]\arrow[l,dashed,bend left=35]&0
\end{tikzcd}\\[2em]
T_p\otimes C)\simeq\mathfrak g\otimes_A C(1)+\check{\mathfrak g}^{*}\otimes_A C
\end{gathered}\right.\quad(**)\)}
Notons que ce dernier isom. est fonctoriel, donc commute à l'action de $A$ sur les deux mbres. Or \inserted{on a vu que} l'action de $A$ sur $\mathfrak g$ est l'action régulière, l'action de $A$ sur $\check{\mathfrak g}^{*}\otimes C$ se fait suivant les hom $A\longrightarrow C$ pour $\sigma\ne1$. Donc la décomposition $(**)$ de $T_p\otimes C$ est celle déduite de la décomposition $(*)$ en \deleted{\unsure{y groupant}} l'écrivant
\[
T_p\otimes C\simeq C_1\oplus\prod_{\sigma\ne1}C_\sigma.
\]
\deleted{Il faut} En tant qu'isomorphisme can., \unsure{est|elle}
\mspage{53}
